\documentclass[11pt,a4paper]{article}
\usepackage[T1]{fontenc}
\usepackage[utf8]{inputenc}
\usepackage{lmodern,microtype}
\usepackage[margin=24mm,headheight=15pt]{geometry}
\usepackage{amsmath,amssymb,amsthm,mathtools}
\usepackage{booktabs,longtable,array,enumitem,xcolor,fancyhdr}
\usepackage{listings}
\usepackage[colorlinks=true,linkcolor=blue!45!black,urlcolor=blue!50!black,citecolor=blue!45!black]{hyperref}
\usepackage{bookmark}
\definecolor{ink}{HTML}{16324F}
\definecolor{muted}{HTML}{546579}
\definecolor{paperblue}{HTML}{EEF3F8}
\hypersetup{pdftitle={DR215: a direct-transfer, reparameterized consistency machine},pdfauthor={Research collaboration for the ProveIt project}}
\pagestyle{fancy}
\fancyhf{}
\fancyhead[L]{\small\textcolor{muted}{DR215 / ProveIt research contribution}}
\fancyhead[R]{\small\textcolor{muted}{21 October 2026}}
\fancyfoot[C]{\thepage}
\renewcommand{\headrulewidth}{0.3pt}
\setlength{\parskip}{0.35em}
\setlength{\parindent}{0pt}
\setlist{itemsep=0.2em,topsep=0.4em}
\newtheorem{theorem}{Theorem}[section]
\newtheorem{lemma}[theorem]{Lemma}
\newtheorem{proposition}[theorem]{Proposition}
\newtheorem{corollary}[theorem]{Corollary}
\theoremstyle{definition}
\newtheorem{definition}[theorem]{Definition}
\newcommand{\N}{\mathbb N}
\newcommand{\Z}{\mathbb Z}
\newcommand{\pair}{\mathsf P}
\newcommand{\imp}{\mathsf I}
\newcommand{\qnt}{\mathsf Q}
\newcommand{\mem}{\mathsf M}
\newcommand{\eqn}{\mathsf E}
\newcommand{\calD}{\mathcal D}
\newcommand{\AX}{\mathsf{AX}}
\newcommand{\MP}{\mathsf{MP}}
\newcommand{\Tr}{\mathsf S}
\newcommand{\Halt}{\operatorname{Halt}}
\newcommand{\Con}{\operatorname{Con}}
\newcommand{\blank}{\mathbf 0}
\newcommand{\file}[1]{\path{#1}}
\newcommand{\code}[1]{\texttt{#1}}
\newcommand{\rowact}[3]{#1\,\mathrm{#2}\,#3}
\numberwithin{equation}{section}
\lstset{basicstyle=\ttfamily\small,backgroundcolor=\color{paperblue},frame=single,rulecolor=\color{paperblue},breaklines=true,columns=fullflexible,keepspaces=true,xleftmargin=0.4em,xrightmargin=0.4em}

\begin{document}
\hypersetup{pageanchor=false}
\begin{titlepage}
\vspace*{1.2cm}
{\color{muted}\large COMPUTABILITY / BUSY BEAVER / ZFC}\par
\vspace{0.7cm}
{\color{ink}\Huge\bfseries DR215}\par
\vspace{0.35cm}
{\color{ink}\LARGE A direct-transfer, reparameterized\par consistency machine}\par
\vspace{0.7cm}
{\Large A concrete 215-state candidate, two new local simulations,\par and the remaining 214-state problem}\par
\vspace{0.8cm}
{\large Prepared for Vladimir Reshetnikov's ProveIt project\par 21 October 2026}\par
\vspace{1.2cm}
\begin{abstract}
We give an exact binary, one-tape Turing machine DR215 with 215 working
states and an ordinary mathematical proof, formalizable in ZFC, that its
blank-tape nonhalting is equivalent to $\Con(\mathrm{ZFC})$. The construction
combines the accepted 35-occurrence single-guard transfer implementation
with a changed sixth-axiom sampler. Its source has 201 distinct subtrees
and 4,732 primitive occurrences. Only two transition rows differ from
the accepted direct-transfer machine, but the axiom produced at a fixed
command argument changes; a corresponding proof compiler is supplied.

Fresh unbounded runtime verification covers all 6,151 guarded primitive
cases and all 430 transition rows. A complete 70,506-transition blank
prefix matches a 177-step synthetic axiom prefix, yielding the protected
startup invariant. Four additional 215-state tables realize two new
predecessor-marker transformations and their compositions; their halting
equivalence holds on every initial tape from state zero.

On the formalism side, antecedent exchange gives a calibrated family of
numerical calculi. An exhaustive 279,936-recipe source census for one
changed row attains 204 source nodes, so that experiment does not improve
the state bound. All positive artifacts, exact tables, finite certificates,
and portable checkers accompany the report. \textbf{No 214-state machine,
new proof-assistant kernel endpoint, or state lower bound is claimed.}
\end{abstract}
\vfill
{\small\color{muted}
Research snapshot: \texttt{b6fb232916d008383e7cf2990d68ddb30fded7db}.\par
Underlying numerical calculus, compiler, and machine architecture are
inherited and explicitly attributed. The new contribution is the checked
composition and supplementary transformations, not a new calibration of
ZFC from first principles.}
\end{titlepage}

\pagenumbering{roman}
\hypersetup{pageanchor=true}
\begingroup
\small
\setlength{\parskip}{0pt}
\tableofcontents
\endgroup
\newpage
\pagenumbering{arabic}

\section{Result and exact scope}
\label{sec:result}

The concrete deliverable is \textbf{DR215}, a deterministic machine with
215 working states, two tape symbols, and a bi-infinite tape. Blank is
zero. The head and start state are zero, and entering a separate halt
label terminates execution. The halt label is excluded from the working
state count. Every working state has exactly two transition rows. The
native table uses historical labels; the dense table has working labels
$0,\ldots,214$ and halt label $215$.

\begin{theorem}[Main result]
\label{thm:main}
For the exact DR215 table supplied with this report,
\[
 \mathrm{ZFC}\ \vdash\quad
 \neg\Halt(\mathrm{DR215},\blank)
 \ \longleftrightarrow\ \Con(\mathrm{ZFC}).
\]
The proof uses the repository's established numerical-calculus-to-ZFC
calibration and arithmetic macro contracts. Its new source, physical
simulation, and startup obligations are identified and checked below.
\end{theorem}

This is the user's 215-state alternative outcome. The 214-state target
remains open in this contribution. The exact primary candidate has the
following measured properties; symbolic rule applications are checker
operations, not transitions of the Turing machine.

\begin{center}
\begin{tabular}{lr}
\toprule
Quantity & DR215\\
\midrule
Working states / transition rows & 215 / 430\\
Primitive source occurrences & 4,732\\
Distinct source subtrees & 201\\
Guarded primitive cases & 6,151\\
Symbolic rule applications & 1,618,678\\
Literal transition rows covered & 430 / 430\\
Blank startup transitions & 70,506\\
Matching synthetic source prefix & 177\\
Explicit startup theorem heads & 10\\
Conservative credit needed at cutoff 1 & 2\\
\bottomrule
\end{tabular}
\end{center}

The new native table has SHA-256
\begin{center}\small
\texttt{fac930e9bc89fed41b631de2f12ef4e48e1c8625ef31ef093d4f14b340f113ad}.
\end{center}
The manifest binds the exact source, dense table, state bijection, checker
code, and verification records. A source DAG count is not a physical state
count: contextual decoding, arithmetic roles, and startup impose additional
constraints.

\subsection{What is new, and what is inherited}

The direct transfer, its contextual source reassociation, the sixth-row
parameter symmetry, and the general predecessor-marker idea all exist in
the inspected repository. Their prior authorship is retained. The new
positive results here are:

\begin{enumerate}
\item The exact composition of the single-guard direct transfer with the
sixth-row sampler change, with its own all-natural physical verification
and its own source and startup binding. It changes the actual axiom
enumeration at fixed command arguments.
\item Two additional marker-role choices, $127\mapsto124$ and
$190\mapsto160$, their independently checked all-tape simulations, and
their compositions with each other and with the earlier $176\mapsto52$
transformation.
\item A calibrated antecedent-exchange family, an exact finite source
census, and bounded source and physical experiments that isolate specific
obstacles to a one-state reduction.
\end{enumerate}

DR215 and the four supplied marker tables are distinct under rooted state
renaming and global left/right reflection from all 84 distinct rooted
graphs represented by 321 parseable 215-state table occurrences in the
frozen comparison corpus. The corpus and its original paths are included.
This proves novelty relative to that corpus, not global discovery priority
and not behavioral inequivalence. The marker machines are deliberately
behaviorally equivalent to their parent.

A final repository check at commit
\texttt{a10b2c302cad2b3f3d0a4ce51022eb1ed5995acb} found the newly
added ZC215 archive. Its main table was separately compared against all
five supplied tables, again with no rooted or reflected match. The late
comparison inputs and receipt are retained separately, so the earlier
84-graph corpus remains exactly reproducible.

\subsection{Verification boundary}

The repository's primary K215 has its own proof-assistant work. That
status is not transferred to these different literal tables. DR215 has
an ordinary proof and executable symbolic induction checks; no new exact
Lean or Rocq endpoint was compiled for it. The checker uses Python and
explicit mathematical induction principles. A complete formal endpoint
would additionally formalize or import that checker soundness and bind
the new literal table, source, and startup certificate.

The whole result also depends on the inherited numerical-calculus
calibration. Checking a propositional recovery certificate or a runtime
certificate alone does not prove the entire equivalence with ZFC. Those
dependencies are stated as named reused theorems, not silently replaced
by tests.

\section{Numerical calculus and the changed sampler}
\label{sec:calculus}

\subsection{Literal arithmetic constructors}

For $a,b\in\N$, define
\[
 \pair(a,b)=\frac{(a+b)(a+b+1)}2+b,
 \qquad \imp(a,b)=4\pair(a,b),
 \qquad \qnt_a b=4\pair(a,b)+2.
\]
Also put $\mem(a,b)=4\pair(a,b)+1$,
$\eqn(a,b)=4\pair(a,b)+3$, falsity $F=1$, truth $0$, and
$NX=\imp(X,1)$. Write $X\to Y=\imp(X,Y)$ and
$\exists_a X=N\qnt_a NX$.

All parameters and theorem codes are arbitrary natural numbers. There
is no well-formedness test, support calculation, or formula parser in
the machine. In particular,
\[
 \imp(0,0)=0,\qquad\operatorname{unpair}(0)=(0,0).
\]
Proofs below use these as literal equalities and do not assume syntactic
injectivity for expressions constructed from arbitrary numerical codes.

The inherited calculus $\calD$ contains truth, numerical modus ponens,
theorem-only generalization, and the following schema families:
\begin{align*}
 R_1(a,b)&=\exists_a\eqn(a,b),\\
 R_2(a,b,c)&=\qnt_a(b\to c)\to(\qnt_a b\to\qnt_a c),\\
 R_3(a,b,c)&=\eqn(a,b)\to(\mem(c,a)\to\mem(c,b)),\\
 R_4(a,b,c)&=\qnt_a\qnt_b c\to\qnt_b\qnt_a c,\\
 R_5(a,b,c)&=\eqn(c,a)\to(\mem(c,b)\to\mem(a,b)),\\
 R_6(a,b,c)&=(\exists_a\qnt_a c)\to c,\\
 R_7(a,b,c)&=\eqn(c,a)\to
       (\eqn(c,b)\to\qnt_{\mem(a,b)}\eqn(a,b)),\\
 L(A,B,C,U)&=((C\to U)\to A)\to((A\to C)\to(B\to C)),
       \qquad U\le A.
\end{align*}

For completeness, the literal gated bundle retained by this construction
is defined by
\begin{align*}
 B_0(b,k,m)&=\qnt_b(m\to N\qnt_b(m\to N\mem(b,k))),\\
 S_0(a,b,k)&=\qnt_k(\mem(k,a)\to\mem(k,b)),\\
 C_0(a,b,k,m)&=\qnt_b\bigl(\mem(b,k)\to
       N\qnt_a(m\to N\qnt_a(S_0(a,b,k)\to B_0(b,k,m)))\bigr),\\
 J_0(a,b,c)&=\qnt_c(C_0(a,b,c,c)\to NB_0(b,c,c)),\\
 V_0(a,b,c)&=\mem(a,b)\to\qnt_c\mem(a,b),\\
 \operatorname{Bundle}(a,b,c)&=N(V_0(a,b,c)\to J_0(a,b,c)).
\end{align*}
The bundle is enabled when $c\ne a$ and $c\ne b$; the generator returns
truth otherwise. These literal gates and the couplings among their
parameters are unchanged.

\subsection{The inherited calibration}

The retained dominated $L$ schema has the same numerical theorem set as
its unrestricted predecessor. The repository supplies a finite recovery
certificate with 3,728 proof nodes and 5,160 term nodes; the candidate
packet retains its checker and certificate. Its side conditions use
$U=0$ or a literal subterm relation that implies $U\le A$ under every
natural substitution. This is a finite metaproof of schema recovery,
not an extra substitution rule in the machine language.

The larger inherited calibration is
\begin{equation}
 \calD\vdash1
 \quad\longleftrightarrow\quad
 \exists p\ \operatorname{Proof}_{\mathrm{ZFC}}(p,\bot).
 \label{eq:calibration}
\end{equation}
We reuse this theorem from the accepted K215 foundation and its
predecessors \cite{k215,direct,rho6,dominated}. No weakening of the set
theory bundle or replacement of ZFC by a potentially weaker theory occurs
in DR215.

\subsection{Changing the family enumerator}

The new sixth generator is
\begin{equation}
 R'_6(a,b,c)=(\exists_a\qnt_a b)\to b.
 \label{eq:rho6new}
\end{equation}
Its third parameter is unused. Every other family remains literally
unchanged. For fixed $a$, the old generator ranges over every matrix
code via $c$, and the new one ranges over every matrix code via $b$.
Consequently their sets of generated natural numbers are equal.

\begin{lemma}[Exact calculus preservation]
The calculus obtained by replacing $R_6$ with $R'_6$ has exactly the
same finite derivability predicate on all natural conclusion codes as
$\calD$. Refutations translate primitively recursively in both directions.
\end{lemma}
\begin{proof}
For an old instance with parameters $(a,b,c)$, use new parameters
$(a,c,a)$, yielding the identical natural code. For a new instance with
parameters $(a,b,c)$, use old parameters $(a,a,b)$. All other proof leaves
and inference steps are unchanged. Replace each sixth-row leaf in a
finite derivation by its identical-code counterpart. There is no
restriction on any matrix code, so this argument also covers junk codes
and zero aliases.
\end{proof}

This does \emph{not} say that an individual AX command has the same
output on the same input tuple. For example, whenever the two matrix
arguments differ, the emitted sixth theorem may differ. That operational
change requires the revised row-selection compiler in
Section~\ref{sec:compiler}. Equality of schema families alone would not
justify silently using the old command compiler.

\section{The exact command language}
\label{sec:commands}

A command store is $(t,w,d,h)$: top, encoded stack, additive counter,
and cached formula. Write $[x;w]=\pair(x,w)$; several heads are nested
to the right over a numerical tail. We use an \emph{explicit} list of
heads when reasoning about stack depth, since zero pairing aliases mean
that numerical list length is not unique.

If $\operatorname{unpair}(w)=(x,v)$, the direct transfer is
\begin{equation}
 \Tr(t,w,d,h)=
 \begin{cases}
 \mathsf{halt},&h=1,\\
 (h,v,d,x),&h\ne1.
 \end{cases}
 \label{eq:transfer}
\end{equation}
It tests the incoming cache only. The older composite
DROP;NIP;NIP also tested the incoming top; its nonhalting store effect
was the same. We use the direct single-guard contract throughout.

For $y=d+x$,
\begin{equation}
 \MP(x)(t,w,d,h)=
 \begin{cases}
 (y,[0;w],y,0),&\imp(t,y)=h,\\
 (0,w,y,0),&\imp(t,y)\ne h.
 \end{cases}
 \label{eq:mp}
\end{equation}
For $\AX(c;x,y)$, set $U=d$, $A=d+x$, $B=h+y$, $C=c$.
Let $G(A,B,C)$ be the enabled bundle or zero. The new AX result is
\begin{equation}
 \begin{split}
 \bigl(G(A,B,C),[&R_1,R_2,R_3,R_4,R_5,R'_6,R_7,\\
                 &L(A,B,C,U),\qnt_U h,t;w],0,0\bigr).
 \end{split}
 \label{eq:ax}
\end{equation}
The $h$ in $\qnt_U h$ is the incoming cache, before adding $y$.
The generated logical rows and gated bundle have implication tag zero;
the generalization has tag two. In particular, none of those freshly
constructed codes is $1$. The tenth stack head saves the incoming top.

The field encodings are
\[
 f(\Tr)=[0],\qquad f(\MP(x))=[1,x],\qquad
 f(\AX(c;x,y))=[c+2,y,x].
\]
The retained source has logical registers
$(N,T,W,A,D,H,X,Y,L)$: next candidate, top, stack, selector, counter,
cache, two scratches, and residual command pointer. At a complete
command boundary, $A=X=Y=0$.

\subsection{The 35-occurrence transfer body}

The direct-transfer source clears $T$ in two primitive occurrences,
copies $H$ to $T$ while restoring $H$ in seven occurrences, then runs
the inherited 26-occurrence NIP routine. The total is $2+7+26=35$.
NIP tests $H=1$; on a nonhalting branch it unpairs $W$ into the new
cache and stack. Thus the full body realizes \eqref{eq:transfer},
preserves $N,L,D$, and returns zero selector and scratches.

All loops decrease natural registers, and the all-natural NIP contract
includes $W=0$. This body and its correct contextual association are
inherited from ST215 \cite{direct}. They are not a new discovery of this
packet.

\subsection{Exact source and table edits}

The pinned direct-transfer source is altered only at four zero-based
primitive ordinals:
\begin{center}
\begin{tabular}{rcc}
\toprule
Ordinal & Old operand & New operand\\
\midrule
958 & register 3 & register 5\\
963 & register 3 & register 5\\
1136 & register 3 & register 5\\
1141 & register 3 & register 5\\
\bottomrule
\end{tabular}
\end{center}
These are the two decrement-and-restore copies of the sixth row's
matrix. Register 3 holds $C$ and register 5 holds $B$ in that AX
context. Changing both references in each copy preserves the copied
register and changes the selected matrix exactly as in
\eqref{eq:rho6new}. Every opcode, source address, and guarded successor
ordinal remains unchanged. The source remains at 201 distinct subtrees,
4,732 occurrences, depth 44, and maximum zero run 6.

The complete native transition delta from ST215 is
\begin{center}
\begin{tabular}{ccc}
\toprule
State/read & Old action & New action\\
\midrule
$199/1$ & $\rowact{1}{R}{166}$ & $\rowact{1}{R}{12}$\\
$206/0$ & $\rowact{0}{R}{186}$ & $\rowact{0}{R}{76}$\\
\bottomrule
\end{tabular}
\end{center}
No state is added or removed. The exact reconstruction uses pinned
parent files and produces both native and dense tables; it does not
invoke a solver. The source is checked using two independent control-flow
resolvers. These finite checks establish literal identities; the physical
runtime check below establishes their unbounded operational meaning.

\section{Unbounded physical simulation}
\label{sec:physical}

A standard tape frame has dispatcher head at integer position $p$,
a marker at $p-1$, source-address bits beginning at $p$, then a zero
gap to a header at $p+64+g$, where $g\in\N$. Register blocks are
\[
 B(v)=0\,1^{v+1}.
\]
Logical registers map to physical block indices
\[
 (N,T,W,A,D,H,X,Y,L)\longmapsto(6,2,0,4,3,7,8,5,1).
\]
After nine primary blocks is a nonempty finite spare list, whose values
are arbitrary naturals. The rest of the tape is zero. The represented
frame includes the entire tape, not only decoded register values.

For every primitive occurrence, decrements split into zero and positive
cases. The supplied runtime derivation has 6,151 cases and quantifies
over arbitrary primary register values, gap, first spare, and remaining
spare list. It uses three kinds of sound inference:

\begin{enumerate}
\item A literal transition from the exact table, with its precise write,
head motion, control target, and clock increment.
\item Natural-number induction over a unary run scanned by a verified
directional self-loop.
\item Structural induction over an arbitrary block list, after deriving
the required one-block action with opaque left and right contexts.
\end{enumerate}

The four block schemas scan $B$ right, convert $B$ to $C=1\,0\,1^v$
right, scan $B$ left, and convert $C$ back to $B$ left. Each has duration
$v+2$. List normalization uses the inductively justified identity
$C(A)\,1=1\,B(A)$. The checker retains all unaffected tape contexts.

Affine expressions have a constant coordinate and twelve nonnegative
variable coordinates: gap, nine primary values, first spare, and width
of the remaining spare list. Each successful case has a finite positive
duration. A positive constant and nonnegative variable coefficients
prevent an infinite source computation from being simulated in finite
binary time. Increment adds a zero spare; no primitive deletes the last
spare, so the frame class is closed.

For DR215, the fresh run completed all 6,151 cases with no failures,
1,618,678 symbolic rule applications, and coverage of all 430 literal
rows. The exact input identities and per-case records are supplied.
This is an unbounded symbolic induction argument, not a test suite over
small integers. It nevertheless remains an ordinary checker-backed
proof with the explicitly stated implementation boundary.

\begin{proposition}[Physical correspondence]
Each source primitive from a standard frame is simulated by a finite
positive number of DR215 transitions, returning the complete frame for
its literal guarded source successor, or entering the separate halt
label exactly for the halting primitive. Consequently finite halting
and infinite source execution lift to the corresponding binary behavior.
\end{proposition}
\begin{proof}
Apply the exact-table symbolic derivation to the appropriate occurrence
and guard, then induct over source steps. Frame closure justifies each
successive application. Positivity of durations supplies divergence
preservation. The unique halt case is checked against the literal
halting row.
\end{proof}

\section{Exact startup and why the opaque tail is safe}
\label{sec:startup}

The complete native blank execution reaches the following frame after
70,506 transitions:
\[
\begin{gathered}
 q=84,\qquad p=-2,\qquad\mathrm{header}=62,\qquad g=0,\\
 \mathrm{PC}=00110010010110010,\qquad\mathrm{ordinal}=58,\\
 (N,T,W,A,D,H,X,Y,L)=(1,0,0,0,0,0,0,18,0).
\end{gathered}
\]
There are 63 zero-valued spare blocks and exactly 100 one-cells. The
certificate records the complete one-cell set; every other tape cell is
zero. Dense-table replay uses the supplied bijection and agrees with the
native run. The altered transition rows do not occur in this finite
startup prefix, but startup is still checked on the new table.

This is an interior AX frame, not a clean command-boundary store.
Starting the source at AX entry ordinal 25 with
\[
 (N,T,W,A,D,H,X,Y,L)=(1,0,2,0,0,0,0,0,0)
\]
reaches the same ordinal and the entire same logical store after exactly
177 primitive steps. Both source-flow resolvers agree. The input is
$\AX(0;0,0)$ on $(t,w,d,h)=(0,2,0,0)$; at those parameters the old and
new sixth-row matrices coincide.

The all-natural AX contract therefore applies to the remaining suffix.
It terminates at a command boundary with $N=1$, $L=A=X=Y=0$,
$t=d=h=0$, and
\[
 w=[R_1(0,0,0),\ldots,R'_6(0,0,0),R_7(0,0,0),
       L(0,0,0,0),\qnt_0 0,0;2].
\]
These ten explicit heads are theorems. The bundle gate is zero at the
all-zero arguments. No literal simulation of the enormous completed
formula codes is needed: exact matching of a finite prefix, terminating
AX semantics, and the universal primitive simulation determine the suffix.

The numerical tail $2$ is not harmless on its own:
\[
 2=\pair(0,1),\qquad1=\pair(1,0).
\]
Repeated unprotected pops would expose falsity. The proof below shows
that the actual schedule never reaches this tail. Historical bootstrap
examples in inherited documents that use ordinal 55 or tail 4 belong to
different tables; the displayed frame and prefix are the DR215 bindings.

\section{The revised strong proof compiler}
\label{sec:compiler}

For every numerical theorem $X$, we construct a finite block $B_X$ with
the following contract for \emph{all} natural $t,w$:
\[
 (t,w,0,0)\xrightarrow{B_X}\mathsf{halt}
 \quad\hbox{or}\quad(X,w,0,0).
 \qquad\text{(Strong)}
\]
The stack need not contain theorems. Early halting is allowed for this
completeness contract; the soundness argument separately establishes
theorem-valued actual execution.

\subsection{Stack operations without an extra opcode}

Put $D=\AX(0;0,0)$ and define
\[
 H=D;\Tr^{11},\qquad
 \operatorname{PUSH}=D;\Tr^9;\MP(0).
\]
After ten transfers following $D$, the store is
$(\qnt_d h,w,0,t)$. The first ten guards see only zero or fresh nonfalsity
heads. The eleventh transfer halts exactly when the original $t=1$;
otherwise it preserves $t$, pops one original stack head into cache,
and resets the counter. Hence for $w=[x;v]$,
\[
 H(t,w,d,h)=\mathsf{halt}\ (t=1),\quad
 H(t,w,d,h)=(t,v,0,x)\ (t\ne1).
\]
In PUSH, nine transfers instead leave the generalized cache, of tag two,
as the tested MP cache. The implication comparison has tag zero, so MP
must fail. Therefore PUSH always returns $(0,[t;w],0,0)$ and never halts.

For zero incoming counter and cache, define
\[
 \operatorname{DISCARD}=\MP(1);\AX(1;0,0);\Tr^{10};\MP(0).
\]
The first MP fails, yields counter one and top zero, and preserves the
stack. AX has $A=C=1$, hence gate zero. Ten transfers expose a quantifier
code while preserving the old stack and zero cache. The last MP fails
and returns $(0,w,0,0)$. This gives the truth block $B_0$.

When top, counter and cache are zero, let $J(0)$ be empty and let
$J(u)=\MP(u)$ for $u>0$. It sets the counter to $u$ without disturbing
the stack. The empty zero case is essential: $\MP(0)$ at all-zero
comparison data would succeed by $\imp(0,0)=0$ and push another zero.

\subsection{The repaired sixth-row selector}

For an unchanged retained row $R_j(a,b,c)$, use
\[
 \operatorname{DISCARD};\AX(c;a,b);\Tr^{j+1};H^{9-j}.
\]
The transfers place the selected row on top. Repeated $H$ then removes
the remaining generated entries and saved zero while retaining that
nonfalsity top; the original stack and zero counter/cache are restored.

For a desired old sixth-row theorem with binder $a$ and matrix $m$,
the new selector is instead
\begin{equation}
 \boxed{\operatorname{DISCARD};\AX(a;a,m);\Tr^7;H^3.}
 \label{eq:newselector}
\end{equation}
After DISCARD, the AX arguments give $A=a$, $B=m$, $C=a$, $U=0$.
The sixth head is exactly $(\exists_a\qnt_a m)\to m$ and its gate is
zero. The same stack-count calculation proves (Strong). This explicitly
repairs the operational effect of changing the sampler.

For $L(a,b,c,u)$ with $u\le a$, use
\[
 \operatorname{DISCARD};J(u);\AX(c;a-u,b);\Tr^9;H.
\]
For an enabled bundle, use
$\operatorname{DISCARD};\AX(c;a,b);H^{10}$. These formulas preserve
the original gates and dominance constraint.

\subsection{Generalization and modus ponens}

Given a compiled theorem $X$, generalization is implemented by
\[
 B_{\qnt_i X}=B_X;\operatorname{PUSH};J(i);\Tr;
                    \AX(i;0,0);\Tr^{10}.
\]
The single transfer loads $X$ into cache without testing the newly
loaded value. AX then puts $\qnt_i X$ in its ninth head, has zero gate
because $A=C=i$, and saves zero. Ten transfers restore the original
stack. This calculation also covers $i=0$ and $X=1$.

For major premise $M=\imp(A,B)$, build the major before the minor:
\[
\begin{split}
 B_B={}&B_M;\operatorname{PUSH};B_A;\operatorname{PUSH};\\
       &\Tr;\Tr;\MP(B);\AX(B;0,0);\Tr^{11}.
\end{split}
\]
After the first four blocks the store is $(0,[A,M;w],0,0)$ unless an
allowed halt has occurred. Two transfers give $(A,w,0,M)$, or halt
already if $A=1$. MP succeeds by the literal equation $M=\imp(A,B)$,
giving $(B,[0;w],B,0)$. The following AX and eleven transfers remove
exactly that additional zero and restore $(B,w,0,0)$, or halt when
$B=1$. For $A=B=0$, MP intentionally succeeds at the native zero alias;
the same final transfer removes its extra stack entry.

Induction on a finite derivation proves (Strong), with the sixth-row
case handled by \eqref{eq:newselector}. A finite proof DAG can be
expanded to a tree; major-before-minor order introduces no logical
dependency. In particular, if $\calD\vdash1$, then $B_1;H$ halts from
every store $(t,w,0,0)$.

\subsection{Cofinal successful candidate numbers}

Let $F(fs,r)$ right-fold a field list by pairing over $r$, and define
\[
 \operatorname{tail}([])=0,\qquad
 \operatorname{tail}(k::L)=1+F(f(k),\operatorname{tail}(L)).
\]
The successor belongs to continuation encoding because a positive
residual is decremented before its next command. A fresh candidate is
decoded directly, without that outer successor.

Set $R=\operatorname{tail}(B_1;H)$. For every $Y\in\N$, the fresh
candidate
\begin{equation}
 N_Y=\pair(2,\pair(Y,\pair(1,R)))
 \label{eq:cofinal}
\end{equation}
starts with $\AX(0;1,Y)$, resetting counter and cache from any retained
store. Its suffix therefore halts by (Strong). Since $N_Y\ge Y$,
successful candidates are unbounded.

Every positive residual strictly decreases after command decoding:
it becomes a right projection of its predecessor after the appropriate
number of unpairings. Each arithmetic command terminates on natural
registers. Thus each nonhalting candidate drains, and the unchanged
enumerator, which increments the candidate number by one, reaches a
chosen successful $N_Y$ unless it has already halted.

\section{Protected credit for the real schedule}
\label{sec:credit}

The following argument uses the actual field arities and candidate
order. It is independent of the values of the axiom matrices, so the
sampler change does not alter it. Let $f$ be the right Cantor projection,
$T_s=s(s+1)/2$, and $W(n)$ the syntactic command word of fresh candidate
$n$, ignoring possible semantic halts. Put $R(0)=[]$ and
$R(r)=W(r-1)$ for positive $r$. If $\operatorname{unpair}(n)=(a,y)$,
\begin{equation}
 W(n)=
 \begin{cases}
 \Tr;R(y),&a=0,\\
 \MP;R(f(y)),&a=1,\\
 \AX;R(f(f(y))),&a\ge2.
 \end{cases}
 \label{eq:schedule}
\end{equation}

Assign conservative weights $\Tr=-1$, $\MP=0$, $\AX=10$; successful
MP's extra truth head is omitted from the lower bound. Let $b(n)$ be
the word's total weight, $B(k)=\sum_{n<k}b(n)$, $g(0)=0$, and
$g(r)=b(r-1)$ for $r>0$. Induction gives $|W(n)|\le n+1$, so every
prefix of $R(r)$ loses at most $r$ units.

\begin{lemma}[Prefix pullback]
If a sequence has nonnegative initial partial sums, composing it with
$f$ or with $f^2$ preserves that property.
\end{lemma}
\begin{proof}
In argument order, the right Cantor projection concatenates
$[0],[0,1],[0,1,2],\ldots$. Every finite prefix consists of complete
lists and an initial part of the next list. Each corresponding sum is
nonnegative. Apply the same argument again for $f^2$.
\end{proof}

\begin{lemma}[Credit bound at cutoff one]
The minimum cumulative conservative weight over every command prefix
of $W(1)W(2)W(3)\cdots$ is $-2$.
\end{lemma}
\begin{proof}
The initial words are $W(0)=\Tr$, $W(1)=\MP$, $W(2)=\Tr\Tr$,
so $B(0),\ldots,B(3)=0,-1,-1,-3$. Jointly induct on Cantor diagonals
to show $B(k)+k\ge0$ and that each diagonal $s\ge2$ is prefix-safe
relative to its start.

Use $q(r)=g(r)+1$. For $k>0$ its prefix sum is
$k+B(k-1)\ge0$. On diagonal $s\ge2$, the first $s-1$ candidates
start with AX, followed by MP and then S. The twice-pulled-back
inequality gives at least $9k$ weight from the first $k$ complete AX
candidates. Before AX number $y$ the balance is at least $9y$;
after its initial $+10$, its continuation loses at most
$f^2(y)\le y$. Every AX prefix is therefore safe.

After the AX portion the balance is at least $9(s-1)$. The MP suffix
can lose at most $s-1$ and the final S word at most $1+s$. Hence
prefixes through the last two words are bounded below by
\[
 9(s-1)-(s-1)-1-s=7s-9>0.
\]
All requested pullback indices lie on earlier diagonals. Combining
diagonal safety with $B(T_s)\ge-3$ and $k\ge3$ preserves
$B(k)+k\ge0$, closing the joint induction.

Starting at cutoff one, the exceptional prefix is
$\MP\Tr\Tr$, of minimum weight $-2$, and no later diagonal reduces
the balance. This establishes both the upper need and its attainment
for the conservative syntactic weights.
\end{proof}

Maintain an explicit finite sequence of theorem heads above a fixed
opaque tail $J$. AX adds ten theorem heads, successful MP adds truth,
failed MP adds none, and S consumes one head. Top and cache remain
theorems whenever no halt has occurred. Starting with ten heads and
debt at most two leaves at least eight explicit heads above $J$ at
every boundary. Before a transfer there is therefore a head to consume;
the tail is never exposed. This uses only exact inversion of
$\pair(x,v)$ and does not assume unique encoded list length.

DR215 starts ordinary enumeration at candidate one with those ten heads
over $J=2$. Its sole halting test can therefore succeed only when the
cache is a theorem equal to $1$.

\section{Proof of the consistency equivalence}
\label{sec:mainproof}

\begin{proof}[Proof of Theorem~\ref{thm:main}]
The exact blank prefix reaches the displayed interior AX frame.
Its matching source prefix and terminating AX contract complete it to
ten theorem heads over tail $2$, with next candidate one and empty
residual. Universal physical correspondence transfers this conclusion
to the literal binary table.

If DR215 halts, the protected-credit invariant and command equations
show that its cache is the numerical theorem $1$. Exact calculus
preservation and \eqref{eq:calibration} then give a finite ZFC refutation.

Conversely, a finite ZFC refutation yields $\calD\vdash1$ by the
inherited calibration. The revised compiler produces the universal
halting word $B_1;H$. The cofinal shell \eqref{eq:cofinal} gives a
successful fresh candidate above the actual cutoff, regardless of the
retained stack. Every preceding nonhalting candidate terminates, so the
enumerator reaches it or halts earlier. Its finite source halting lifts
to finite binary halting.

The resulting equivalence
\[
 \Halt(\mathrm{DR215},\blank)
 \longleftrightarrow
 \exists p\ \operatorname{Proof}_{\mathrm{ZFC}}(p,\bot)
\]
uses finite calculations, primitive-recursive translations, and
induction on natural numbers, finite lists, and finite derivations.
These arguments are formalizable in ZFC without assuming consistency.
Negating the equivalence proves the theorem.
\end{proof}

No unconditional nonhalting claim is made. If ZFC is consistent, the
machine does not halt and the usual second-incompleteness argument
prevents ZFC from proving this nonhalting statement. Mere consistency is
not being used to assert that ZFC cannot prove a false halting statement.

\section{Two new predecessor-marker transformations}
\label{sec:markers}

The supplementary family starts from the primary K215 native table,
not from DR215. Its parent SHA-256 is
\begin{center}\small
\texttt{9683d6ed9628e6ce407c315b08f5da3d063da94e92706f275ca0f33e8272e4b2}.
\end{center}
Each transformation merges two logical boundary roles and reuses the
freed physical label as an internal selector. Consequently the state
count stays at 215. The gain is a different physical organization with
a simple all-tape simulation, potentially useful for later sharing.

\subsection{M124: distinguish old roles 124 and 127}

The old states have a common read-one action $1R138$, while their
read-zero actions are $0R62$ and $0R14$. The only incoming row to
124 is $104/1:1R124$, and the only incoming row to 127 is
$110/0:0R127$. Therefore, on every entry from start zero, the cell
immediately left of the head records the role: one for 124 and zero
for 127, on every initial tape.

Map old role 127 to boundary role 124 and reuse physical state 127 as
the selector. Exactly four rows change:
\begin{center}
\begin{tabular}{ccc}
\toprule
Row & Old & New\\
\midrule
$110/0$ & $0R127$ & $0R124$\\
$124/0$ & $0R62$ & $0L127$\\
$127/0$ & $0R14$ & $0R8$\\
$127/1$ & $1R138$ & $1R59$\\
\bottomrule
\end{tabular}
\end{center}
Provider rows $8/0:0R14$ and $59/0:0R62$ complete the simulated
step. On current zero the new machine goes left to inspect the
unchanged marker, right to the selected provider, then right to the
old target. Both cells retain their original symbols. On current one,
the common old action executes directly.

\subsection{M160: distinguish old roles 160 and 190}

Here the common action is read-zero $0R99$, while the read-one actions
are $1R99$ and $1R52$. The only incoming rows are
$139/0:0R160$ and $165/1:1R190$, so the predecessor marker is zero
for old 160 and one for old 190. Map $190\mapsto160$ and change:
\begin{center}
\begin{tabular}{ccc}
\toprule
Row & Old & New\\
\midrule
$160/1$ & $1R99$ & $1L190$\\
$165/1$ & $1R190$ & $1R160$\\
$190/0$ & $0R99$ & $0R86$\\
$190/1$ & $1R52$ & $1R193$\\
\bottomrule
\end{tabular}
\end{center}
Provider rows $86/1:1R99$ and $193/1:1R52$ give the same exact
three-step simulation, this time on exceptional current symbol one.

\subsection{Composition and global proof}

The M124 and M160 changed rows are disjoint, and neither modifies a
provider row used by the other. They compose into an eight-row edit.
They also compose with the repository's earlier $176\mapsto52$
conditional-marker transformation, producing a fourteen-row edit.
For that inherited transformation the exceptional incoming route
$189/1;156/1$ forces the current cell to one and cannot trigger its
read-zero discriminator; the other incoming routes write the required
zero predecessor. The builder and independent auditor check the entire
incoming-edge enumeration.

\begin{theorem}[All-tape halting equivalence]
For each of the four supplied marker variants, starting old and new
machines in state zero at the same head position on any binary tape,
the machines halt together. There is a checkpoint map with exact whole
tape and head equality, control mapped by the stated role identifications,
and times $t(n)$ satisfying $n\le t(n)\le3n$.
\end{theorem}
\begin{proof}
Induct on old steps. A nonexceptional row is a literal copy with targets
mapped by the role map, so it uses one new step. An exceptional row uses
the verified three-step marker/provider calculation. Incoming-edge
invariants establish the required marker at every such boundary.
The independent check compares both local cells and confirms that the
block reads nothing outside them, leaving arbitrary exterior tape intact.

Set $t(n+1)-t(n)$ to one or three accordingly. Internal selector and
provider states are working states, so no exceptional block halts inside.
A halting transition occurs at corresponding boundaries. The bound
$t(n)\ge n$ rules out a finite new run representing infinitely many old
steps, yielding the converse as well.
\end{proof}

This proves the same ZFC consistency equivalence as the accepted parent.
It is stronger than a blank-tape simulation, but makes no claim about
starting arbitrarily inside a replaced role. Finite startup checks are
corroboration rather than the basis of the all-tape proof: M124 and its
compositions reach the old 95-one-cell seed after 95,840 transitions;
M160 reaches it after 95,714. Every mapped old-step boundary and the
complete final tape were compared.

\section{A calibrated family of changed formalisms}
\label{sec:formalism}

The user's permission to change the formalism is substantial: only
refutability, not theorem-for-theorem identity, ultimately needs to be
preserved. The positive formalism experiments here establish the stronger
identity property, because it supplies a short trustworthy calibration
while allowing a different axiom layout.

\subsection{Antecedent exchange on numerical codes}

Retain the dominated $L$ basis, which recovers the unrestricted basis
and hence the propositional $K$ and $S$ schemata. They derive
\begin{equation}
 (X\to(Y\to Z))\to(Y\to(X\to Z)).
 \label{eq:exchange}
\end{equation}
For an explicit derivation template, assume $f:X\to(Y\to Z)$,
$y:Y$, and $x:X$. Two MP steps yield $Z$. The finite $K/S$ deduction
construction discharges $x$, then $y$, then $f$. It replaces the
discharged assumption by $A\to A$, any other assumption or theorem
by a $K$ instance, and an MP step by an $S$ instance. No GEN is used,
so no generalization side condition is hidden. Exchange $X,Y$ for the
reverse implication.

Choose any subset of the following four rows and exchange their two
curried antecedents:
\begin{center}
\small
\begin{tabular}{clll}
\toprule
Row & $X$ & $Y$ & $Z$\\
\midrule
$R_2$ & $\qnt_A(B\to C)$ & $\qnt_A B$ & $\qnt_A C$\\
$R_3$ & $\eqn(A,B)$ & $\mem(C,A)$ & $\mem(C,B)$\\
$R_5$ & $\eqn(C,A)$ & $\mem(C,B)$ & $\mem(A,B)$\\
$R_7$ & $\eqn(C,A)$ & $\eqn(C,B)$ & $\qnt_{\mem(A,B)}\eqn(A,B)$\\
\bottomrule
\end{tabular}
\end{center}
Both old and changed calculi retain the same $K/S$ basis. Each derives
every missing axiom from its partner by \eqref{eq:exchange} and one MP.
Replacing axiom leaves gives primitive-recursive translations of all
finite proofs, preserving every conclusion code, including truth and
falsity. The argument applies at arbitrary natural parameter values and
requires no formula parser or ordinary syntactic substitution theorem.

Independent permutations of the three parameter positions in a row
also preserve its full schema family. The command compiler selects an
instance by resetting counter/cache and using the inverse permutation
for the AX fields, as the sixth-row example demonstrates. All gates
and private parameter couplings in the set theory bundle remain unchanged.

\subsection{Measured source cost}

In the primary K215 source, the merged equality row $R_7$, called $G$
by the source tools, occupies primitive interval $[479,888)$. Its first
two atom constructors occupy $[479,643)$. Swapping registers four and
five in the CI/CD operands of that prefix exchanges exactly those two
antecedents, without changing any source address, opcode, or guarded
successor. The resulting source has 205 distinct subtrees, against
the primary baseline's 202.

An exhaustive census then applies all $6^7=279{,}936$ independent
parameter permutations of the seven logical rows to that exchanged
source. At the retained literal association, the minimum is 204 source
nodes, achieved by four recipes. All recipes keep 4,751 primitive
occurrences. A subsequent width-64 closed-reassociation beam visited
76,675 sources over eleven rounds and saved no source below 204.
These counts are scoped to those finite or bounded search families.

We also reran both finite families on DR215 itself. All 279,936 row
permutations have minimum 201, with four minimizers. After antecedent
exchange in $R_7$, the same exhaustive census has minimum 203, again
with four minimizers. A width-64 closed-reassociation beam on DR215
visited 75,077 sources over eleven rounds and retained minimum 201.
The portable replays reproduced the histograms and every minimizing
source byte string. Thus no source of 200 or fewer nodes was found in
these precisely specified families.

\begin{center}
\begin{tabular}{lrr}
\toprule
Single changed row / construction & Best saved DAG & Zero run\\
\midrule
Exchanged $R_2$, closed-row re-emission & 205 & 6\\
Exchanged $R_3$, closed-row re-emission & 205 & 6\\
Exchanged $R_5$, closed-row re-emission & 207 & 6\\
Exchanged $R_7$, closed-row re-emission & 206 & 7\\
Exchanged $R_7$, original association & 205 & 6\\
Dual $R_6$, experimental and uncalibrated & 203 & 8\\
\bottomrule
\end{tabular}
\end{center}

The last row experiments with $C\to\qnt_A\exists_A C$ in place of
$(\exists_A\qnt_A C)\to C$. Its apparently familiar classical duality
does not establish the needed result: the usual derivation uses ex falso,
while the removed original row itself supplies ex falso in this numerical
basis. The repository already contains a countermodel to equality of the
two theorem sets. A separate relative refutation translation could still
be useful, but none is proved here. That row is not an accepted candidate
and contributes no machine bound.

\section{The 214-state frontier: concrete obstructions and next targets}
\label{sec:frontier}

\subsection{Why fewer source nodes did not automatically save a state}

The attempted combination of direct transfer with MP-first command order
has a valid abstract compiler and a better credit bound: only one
protected head is needed at cutoff one. Some packed layouts have the
same 201-node source size as ST215. That did not suffice for a physical
215-state allocation in the tested construction.

The old contextual implementation shares the CD1 and CD5 clear roles by
inspecting a predecessor bit. Moving the direct-transfer clear under the
MP-first branch changes that bit from one to zero, requiring an additional
generic CD1 role unless a different contextual implementation is found.
Early allocations passed source-prefix decoding but altered an arithmetic
row used by CD8. Once the actually used $134/0$ row was protected, the
tested 201-node layouts became infeasible under that allocator's fixed-role
policy. These rejected tables are not presented as valid candidates.

The lesson is specific: prefix decoding alone is insufficient, and source
DAG cardinality is only one part of the physical allocation problem.
The positive DR215 packet avoids inheriting a partial allocator mask as a
runtime proof; it rederives all guarded physical cases for its actual table.

\subsection{Fusing the transfer copy with its guard}

A separate source-only improvement reduces the direct transfer body from
35 to 30 primitive occurrences. Clear $T$ once, then pair every successful
decrement of $H$ in the cache-equals-one test and residual cleanup with
an increment of $T$. At the completed pairs, $T+H$ equals the incoming
cache. The zero case skips the halt with $T=H=0$; the one case halts;
the larger cases exhaust $H$ into $T$ and enter the unchanged Cantor
unpair suffix. The latter restores the scratches and yields exactly the
single-guard transfer contract.

Inserting an increment changes the loop's ancestor depth, so its comparison
continue must be CC3 rather than CC2. The initial exploratory source missed
that correction; only the corrected source is supplied. Its complete
program has 4,727 occurrences but 208 distinct subtrees, against 4,732
and 201 in its parent. Thus it saves occurrences while losing sharing.
The addendum includes the literal recipe, an all-natural source proof,
dual-CFG checks and finite corroboration. There is no binary allocation,
physical runtime certificate, or state-count claim for this source.

\subsection{Sharing one selector between two marker mergers}

The two new marker mergers suggest a genuine one-state target. M124
distinguishes current zero and M160 distinguishes current one. If their
two internal selectors could be replaced by one shared selector, one
physical label would become removable.

In the simplest three-step version, the marker-zero provider would need
both actions
\[
 0\mapsto0R14,\qquad1\mapsto1R99,
\]
while the marker-one provider would need
\[
 0\mapsto0R62,\qquad1\mapsto1R52.
\]
No existing unchanged provider state supplies either required pair.
Longer contextual provider blocks and temporary writes enlarge the search
space, but must restore the complete tape and head position before
rejoining the mapped old execution.

The supplied bounded shared-selector search tests both choices of retained
helper label, temporary current-bit encodings, directions and writes of
the two helper rows and two rewritten entry rows, and matching local
contexts. Its radius-seven, fuel-97 primary-table census contains
46,895,104 specifications across the M124/M160 and earlier-M052/M160
families and has no survivors. This is a finite negative result for the
stated local simulation model, not a proof that a shared selector or a
214-state machine is impossible. Longer simulations, different invariants,
more simultaneous edits, or a changed arithmetic representation remain
outside that conclusion.

The broader final census uses DR215 and all 200 disjoint pair combinations
from a recorded inventory of 21 marker pairs: 105 combinations have
opposite exceptional symbols and 95 have the same symbol. Temporary
writes can encode pair identity even in the latter case. Both exceptional
entry rows move left. Each family has
\[
 2\cdot4\cdot(2\cdot2\cdot214)^2=5{,}861{,}888
\]
parameter specifications, counting selector choice, temporary writes, and
each selector row's write, direction and target. Across all families,
1,172,377,600 specifications were covered: 707,406,332 individually
evaluated and 464,971,268 uniformly pruned, with no survivors. Pruning
applies only when a failed path never reads the innermost variable row,
so every remaining assignment to that row has the same failure. These
are specifications, not necessarily distinct machines. Radius seven,
fuel 97, a selected role inventory, an overapproximated context domain,
and the prescribed edit shape limit the conclusion. The included C++17
sources, exact inventories and wrapper reproduce the census.

\subsection{What would count as the next substantive advance}

Three directions remain concrete enough to test:
\begin{enumerate}
\item A shared-selector construction with a proved invariant outside the
present finite-window family. The displayed provider requirements make
the missing interface explicit.
\item A direct MP-first encoding that implements its moved clear routine
without the extra generic role. Its abstract compiler and protected-credit
theorem are already available; its physical binding is the missing part.
\item A genuinely smaller refutation-equivalent numerical basis. Unlike
the positive exchange family, this may change theorem sets, but it needs
both primitive-recursive refutation directions. A classical equivalence
argument that reuses an axiom being removed is not enough.
\end{enumerate}

The measured failures help prioritize these targets. They establish no
minimum-state theorem and do not justify calling 215 optimal. In
particular, we have not supplied a smaller class theory, an arithmetic
replacement for the set theory bundle, or a calibrated dual-sixth-row
formalism; those would be different research projects with additional
proof obligations.

\section{Reproduction and artifact map}
\label{sec:reproduce}

The archive contains this editable \LaTeX{} source, its compiled PDF,
the main candidate packet, marker family, formalism appendix, comparison
corpus, licenses and provenance, and machine-readable hashes. The
top-level README gives the exact final command lines. Positive candidate
verification requires Python 3 with assertions enabled and no network or
external repository checkout. Optional synthesis experiments may have
additional dependencies, stated by their own scripts.

The default candidate replay reconstructs the literal source and both
tables from the pinned parent; checks complete source/control differences;
replays startup and the synthetic AX prefix; checks the dominated-row
certificate; and audits the full symbolic case inventory. Its full-runtime
option freshly rederives all 6,151 induction cases. The marker checker
reconstructs all four tables and checks the finite premises of the
all-tape theorem; a separately written auditor uses no builder imports.
The graph audit compares rooted labeled transition graphs and reflections
against the included inert corpus.

Runtime records and research receipts can retain original scratch paths
as provenance. The positive reproduction commands use included relative
inputs. They do not depend on those historical paths. SHA-256 manifests
identify byte strings; they provide integrity checking, not a substitute
for the mathematical arguments.

\subsection{Attribution}

The underlying architecture, native numerical theory, dominated-row
calibration, register compiler, runtime induction engine, and direct-transfer
and rho6 predecessors belong to the ProveIt contribution history. Its
retained notices credit, within their respective scopes, ProveIt
contributors, Elliot Glazer's compact theory work, Jeff's busy-beaver/Rocq
reference, Andrew James Wade's compiler ancestry, BusyCoq, and the Coq
Library of Undecidability. This contribution's exact composition, marker
instances, exchange family, audits, and report were prepared by the
research collaboration for Vladimir Reshetnikov. Independent agent
reviews are internal review, not an external human endorsement.

New material follows the repository's MIT No Attribution terms; included
third-party components retain their original terms. Original notices and
license files are supplied. No blanket ownership or global priority claim
is made.

\clearpage
\begin{thebibliography}{9}
\small
\raggedright
\bibitem{k215} ProveIt contributors.
\emph{Candidate K215 and accepted calibration}, pinned repository snapshot.
\url{https://github.com/VladimirReshetnikov/ProveIt/tree/b6fb232916d008383e7cf2990d68ddb30fded7db/Computability/BusyBeaver/Research/ZFC/candidate-k215}
\bibitem{direct} ProveIt contributors.
\emph{ST215: direct transfer through contextual reassociation}.
\url{https://github.com/VladimirReshetnikov/ProveIt/blob/b6fb232916d008383e7cf2990d68ddb30fded7db/Computability/BusyBeaver/Research/ZFC/candidate-k215/alternatives/direct-transfer/report.tex}
\bibitem{rho6} ProveIt contributors.
\emph{K215-R6: two fully checked 215-state machines}.
\url{https://github.com/VladimirReshetnikov/ProveIt/blob/b6fb232916d008383e7cf2990d68ddb30fded7db/Computability/BusyBeaver/Research/ZFC/candidate-k215/alternatives/rho6/README.md}
\bibitem{dominated} ProveIt contributors.
\emph{Dominated-rho8 alternative: proof}.
\url{https://github.com/VladimirReshetnikov/ProveIt/blob/b6fb232916d008383e7cf2990d68ddb30fded7db/Computability/BusyBeaver/Research/ZFC/candidate-k218/alternatives/dominated/proof.md}
\bibitem{mpfirst} ProveIt contributors.
\emph{MP-first transfer: protected-credit proof}.
\url{https://github.com/VladimirReshetnikov/ProveIt/blob/b6fb232916d008383e7cf2990d68ddb30fded7db/Computability/BusyBeaver/Research/ZFC/candidate-k215/alternatives/mp-first-transfer/proofs/permuted_credit.md}
\bibitem{conditional} ProveIt contributors.
\emph{Conditional-marker alternative}.
\url{https://github.com/VladimirReshetnikov/ProveIt/tree/b6fb232916d008383e7cf2990d68ddb30fded7db/Computability/BusyBeaver/Research/ZFC/candidate-k215/alternatives/conditional-marker}
\end{thebibliography}

\clearpage
\appendix
\section{Complete DR215 dense transition table}
\label{sec:table}

Each entry $wDz$ means write $w$, move in direction $D$, then enter
state $z$. The four groups of columns merely wrap the state list for
printing. State zero starts on blank tape; $215$ is the separate halt
label. The machine file and explicit native/dense bijection are the
authoritative machine-readable presentations.

\begingroup\fontsize{8.5}{10}\selectfont\ttfamily
\setlength{\tabcolsep}{4pt}
\begin{longtable}{rll@{\hspace{1em}}rll@{\hspace{1em}}rll@{\hspace{1em}}rll}
\toprule
$q$ & $0$ & $1$ & $q$ & $0$ & $1$ & $q$ & $0$ & $1$ & $q$ & $0$ & $1$\\
\midrule
\endfirsthead
\toprule
$q$ & $0$ & $1$ & $q$ & $0$ & $1$ & $q$ & $0$ & $1$ & $q$ & $0$ & $1$\\
\midrule
\endhead
\midrule
\multicolumn{12}{r}{\normalfont\footnotesize Continued on next page}\\
\endfoot
\bottomrule
\endlastfoot
0 & 1R91 & 1R0 & 54 & 0R22 & 1R54 & 108 & 0R148 & 1R183 & 162 & 0L108 & 1R170\\
1 & 0L1 & 1L2 & 55 & 0L11 & 0R50 & 109 & 0L204 & 0L94 & 163 & 0R144 & 1R169\\
2 & 0L3 & 1L2 & 56 & 0R60 & 0L16 & 110 & 0R123 & 1R103 & 164 & 0R193 & 1R28\\
3 & 0L5 & 1L2 & 57 & 0R69 & 1R70 & 111 & 0R206 & 1R180 & 165 & 0R158 & 1R141\\
4 & 1R6 & 0L4 & 58 & 0R61 & 1R71 & 112 & 0R162 & 1R178 & 166 & 0R175 & 1R175\\
5 & 0L9 & 1L2 & 59 & 0R54 & 1R59 & 113 & 0R146 & 1R144 & 167 & 0R178 & 1R121\\
6 & 0R6 & 0R10 & 60 & 0R48 & 0R25 & 114 & 0R157 & 1R197 & 168 & 0R143 & 1R162\\
7 & 0R83 & 1R187 & 61 & 0L28 & 1L72 & 115 & 0L130 & 0L1 & 169 & 0R75 & 1R60\\
8 & 0R14 & 1R13 & 62 & 0L116 & 1L30 & 116 & 0R38 & 1R72 & 170 & 0R143 & 1R87\\
9 & 0L15 & 1L2 & 63 & 0R73 & 1R63 & 117 & 0R61 & 1R131 & 171 & 0R103 & 1R95\\
10 & 0L16 & 0R10 & 64 & 0R96 & 1R148 & 118 & 0R61 & 1R132 & 172 & 0R138 & 1R83\\
11 & 0L26 & 1R128 & 65 & 1R181 & 1L53 & 119 & 0R70 & 0L16 & 173 & 0R172 & 1R75\\
12 & 0R19 & 1R20 & 66 & 0R142 & 1R143 & 120 & 0R14 & 1R131 & 174 & 0R20 & 1R37\\
13 & 0L88 & 1L31 & 67 & 0R76 & 1R215 & 121 & 0R21 & 1R133 & 175 & 0R140 & 1L147\\
14 & 0L11 & 0R22 & 68 & 0L193 & 1R31 & 122 & 0R145 & 1R175 & 176 & 0R96 & 1R51\\
15 & 0L23 & 1L2 & 69 & 0R40 & 1R62 & 123 & 0R135 & 1L20 & 177 & 0R165 & 1R51\\
16 & 1L1 & 0L1 & 70 & 0L88 & 1R77 & 124 & 0R129 & 1R177 & 178 & 0R17 & 1L130\\
17 & 0R24 & 0R32 & 71 & 0R55 & 1L78 & 125 & 0R201 & 1R168 & 179 & 0R76 & 1R115\\
18 & 0R198 & 1R197 & 72 & 0R79 & 1R80 & 126 & 0R75 & 1R173 & 180 & 0R105 & 1R123\\
19 & 0R28 & 1R29 & 73 & 1R0 & 1R63 & 127 & 0R17 & 1R212 & 181 & 0L39 & 1L109\\
20 & 0R30 & 1R31 & 74 & 0R74 & 1R81 & 128 & 0R75 & 1R134 & 182 & 0R168 & 1R12\\
21 & 0R24 & 0R59 & 75 & 0L143 & 1L139 & 129 & 0R12 & 1R139 & 183 & 0R148 & 1R114\\
22 & 0R33 & 1R22 & 76 & 0R84 & 1R85 & 130 & 0R30 & 1R69 & 184 & 0R11 & 1R97\\
23 & 0L34 & 1L2 & 77 & 0R88 & 1L131 & 131 & 0R38 & 1L123 & 185 & 0R189 & 1R98\\
24 & 0L35 & 0R0 & 78 & 0L107 & 1R56 & 132 & 0R151 & 0L16 & 186 & 0R144 & 1R195\\
25 & 0R36 & 1R25 & 79 & 0R48 & 0R33 & 133 & 0R152 & 0L16 & 187 & 0R154 & 1R188\\
26 & 0R127 & 1R37 & 80 & 0R90 & 0L41 & 134 & 0R158 & 1R137 & 188 & 0R75 & 1R182\\
27 & 0R38 & 0L39 & 81 & 0R194 & 1R174 & 135 & 0R24 & 0R47 & 189 & 0R64 & 1R7\\
28 & 0R40 & 1R41 & 82 & 0R95 & 1R12 & 136 & 0R138 & 1R177 & 190 & 0R148 & 1R126\\
29 & 0R42 & 1R43 & 83 & 0R87 & 1R96 & 137 & 0R191 & 1R93 & 191 & 0R153 & 1R111\\
30 & 0R44 & 1R45 & 84 & 0R203 & 1R21 & 138 & 0R95 & 1R156 & 192 & 0R150 & 1R122\\
31 & 0R21 & 1R46 & 85 & 0R115 & 1R99 & 139 & 0R172 & 1R163 & 193 & 0R89 & 1R87\\
32 & 0R47 & 1R32 & 86 & 0R148 & 1R112 & 140 & 0R12 & 1R156 & 194 & 0R214 & 1R115\\
33 & 0R32 & 1R33 & 87 & 0R101 & 1R102 & 141 & 0R87 & 1L163 & 195 & 0R202 & 1R108\\
34 & 0L49 & 1L2 & 88 & 0R79 & 1R78 & 142 & 0R87 & 1R155 & 196 & 0R148 & 1R166\\
35 & 1R35 & 1L1 & 89 & 0R100 & 1R171 & 143 & 0L150 & 1L160 & 197 & 0R196 & 1R157\\
36 & 0R50 & 1R36 & 90 & 0R104 & 1R105 & 144 & 0R83 & 1R87 & 198 & 0R138 & 1R136\\
37 & 0R13 & 1R52 & 91 & 0R65 & 0R0 & 145 & 0R156 & 1R12 & 199 & 0R192 & 1R18\\
38 & 0L11 & 0R36 & 92 & 0R27 & 1R106 & 146 & 0R82 & 1R141 & 200 & 0R87 & 1R205\\
39 & 0L4 & 1L53 & 93 & 0R179 & 0L204 & 147 & 0R161 & 1R165 & 201 & 0R160 & 1R190\\
40 & 0R24 & 0R25 & 94 & 1L109 & 1L94 & 148 & 0L66 & 1L139 & 202 & 0R125 & 1R149\\
41 & 0R13 & 0L16 & 95 & 0R19 & 1R110 & 149 & 0R146 & 1R199 & 203 & 0R44 & 1R99\\
42 & 0R55 & 1R56 & 96 & 0R60 & 1R60 & 150 & 0R156 & 1R113 & 204 & 0L4 & 0L204\\
43 & 0R57 & 1R58 & 97 & 0R164 & 1R154 & 151 & 0R96 & 1R96 & 205 & 0R193 & 1R138\\
44 & 0R24 & 0R54 & 98 & 0R66 & 1R186 & 152 & 0R48 & 0R32 & 206 & 0R161 & 1R207\\
45 & 0R60 & 1R61 & 99 & 0R104 & 0L204 & 153 & 0R68 & 1R200 & 207 & 0R208 & 1R179\\
46 & 0R62 & 0L16 & 100 & 0R116 & 1R117 & 154 & 0L207 & 1R95 & 208 & 0R141 & 1R209\\
47 & 0R25 & 1R47 & 101 & 0R71 & 1R58 & 155 & 0R165 & 1R176 & 209 & 0R142 & 1R210\\
48 & 0L35 & 0R63 & 102 & 0R57 & 1R118 & 156 & 0R19 & 1R167 & 210 & 0R211 & 1R86\\
49 & 0L74 & 1L2 & 103 & 0R21 & 1R119 & 157 & 0R170 & 1R162 & 211 & 0R185 & 1R162\\
50 & 0R16 & 1R50 & 104 & 0R21 & 0L10 & 158 & 0R156 & 1R95 & 212 & 0R184 & 1L72\\
51 & 0R145 & 1R141 & 105 & 0R61 & 1R41 & 159 & 0R124 & 1R161 & 213 & 0R92 & 1R107\\
52 & 0R67 & 1R68 & 106 & 0R120 & 1R8 & 160 & 0R96 & 1R159 & 214 & 0R213 & 1R127\\
53 & 0L39 & 1L53 & 107 & 0R100 & 1R121 & 161 & 0L7 & 1R140 &  &  & \\
\end{longtable}
\endgroup


\end{document}
